\documentclass[10pt,a4paper]{amsart}
\usepackage[margin=19mm]{geometry}
\usepackage{amsmath,amssymb,mathtools}
\usepackage[scr=rsfs,cal=euler]{mathalfa}
\usepackage{xfrac}
\usepackage[unicode,hidelinks,bookmarks=false]{hyperref}
\newcommand{\ie}{i.e.\ }
\newcommand{\cf}{cf.\ }
\newcommand{\resp}{respectively\ }
\newcommand{\Subsection}{\subsection}
\providecommand{\nobreakdash}{\nobreak\hbox{-}}
\setlength{\emergencystretch}{2em}
\multlinegap0pt
\usepackage{amssymb}
\usepackage{bm}
\usepackage{mathtools}
\usepackage{dsfont}
\usepackage{enumerate}
\newtheorem{theorem}{Theorem}
\newtheorem{lemma}{Lemma}
\def\C{{\mathbb C}}
\def\R{{\mathbb R}}
\def\Z{{\mathbb Z}}
\newcommand{\norm}[1]{\left\lVert #1 \right\rVert}
\newcommand{\supp}{{\mathrm{supp}}}
\title{Universal completeness of exponentials}
\author{Susanna Bertolini, Enric Florit-Simon, Lukas Liehr, Mitchell A. Taylor}
\date{Source: arXiv:2609.20805v1 (CC BY 4.0)}
\begin{document}
\maketitle

\noindent\emph{Source and licence.} Universal completeness of exponentials. Version arXiv:2609.20805v1 (CC BY 4.0), distributed by arXiv under the Creative Commons Attribution 4.0 International licence, http://creativecommons.org/licenses/by/4.0/, as recorded in the arXiv licence metadata for this exact version and shown on the abstract page https://arxiv.org/abs/2609.20805v1. Full licence notice: \texttt{licenses/arxiv-2609.20805-CC-BY-4.0.txt}.\par\smallskip

\noindent\textit{Editorial scope.} Counted unit: the article's theorem thm:null (source0.tex lines 277--279). The article's complete original proof of it is delivered with it (source0.tex lines 951--1232, one \texttt{proof} environment of 282 lines, which the article places away from the statement). No mathematics is written, repaired or re-derived: the statement and the proof are the article's own lines, in the article's own notation, and no retained line is substituted or re-flowed.  Retained uncounted as the source-local context the proof needs: lines 865--874; lines 898--901.  Nothing was omitted from any retained span, so the package carries no omission record.  No retained line contains a citation, so the wrapper carries no bibliography and no bibliography entry.  No retained line contains a figure environment, an image-inclusion command or drawing code, so no image is delivered and the package declares no asset dependency.  Editorial formatting only, in addition: an AMS \texttt{amsart} preamble that restates the article's own declarations and macros used by the retained lines, namely the environments \texttt{theorem}, \texttt{lemma} on separate counters, together with the standard packages those lines need.  The retained environments print in this document as Theorem 1, Lemma 1 in order of appearance, whereas the article prints the same items with the numbering of its own full document.  The complete native source of the article as distributed in the arXiv e-print for this exact version is bundled unchanged as \texttt{sources/210916/source0.tex}.
\begin{lemma}
\label{lem:asymp-thin-interpolation}
Let 
\(I\subseteq\Z\) finite, and let
\((v_k)_{k\in I} \subset\C\). Given \(0<\varepsilon\leq1\), there exists
\(f\in L^2(\R)\), with $\supp (f)\subset [0,1]$, $|\supp (f)|\leq \varepsilon$, and $\hat f(\lambda_k)=v_k$ for $k\in I$.
Moreover,
$$\quad
    \norm{f}_2^2\leq\frac{C}{\varepsilon}\sum_{k\in I}|v_k|^2.$$
\end{lemma}

\begin{equation}\label{eq:asymp-bessel}
    \sum_{n\in\Z}|\widehat h(\lambda_n)|^2
    \leq C\|h\|_2^2.
\end{equation}

\begin{theorem}\label{thm:null}
    Let $\Lambda = \{ n+ \delta_n : n \in \Z \}$, with $\delta_n\to 0$ as $|n|\to\infty$. For every $\varepsilon>0$, there exists a measurable $S \subseteq \R$ such that $|S| < \varepsilon$ and $E(\Lambda)$ is not complete in $L^2(S)$.
\end{theorem}

\begin{proof}[Proof of Theorem~\ref{thm:null}]
Fix \(\varepsilon>0\). In view of the preceding reductions, we may
assume that
\[
    0<|\delta_n|<\frac18,\quad n\in\Z.
\]
We divide the construction into three steps.

\medskip
\noindent\textbf{Step 1: Adding zeros on a fixed dyadic block.}
For \(j\geq3\), let
\[
    B_j=\{n\in\Z:2^{-j-1}\leq|\delta_n|<2^{-j}\}.
\]
Each \(B_j\) is finite, since \(\delta_n\to0\), and the blocks
\((B_j)_{j\geq3}\) form a partition of \(\Z\).

For \(j\geq3\), define
\[
    \varepsilon_j
    =\frac{\varepsilon}{4}\,
    \frac{
        2^{-j}+\displaystyle\max_{n\in B_j}\frac1{(1+|n|)^2}
    }{
        \displaystyle\sum_{r\geq3}2^{-r}
        +\displaystyle\sum_{n\in\Z}\frac1{(1+|n|)^2}
    },
\]
where the maximum is understood to be zero when \(B_j=\varnothing\). These numbers are chosen so that, for some constant
\(C_\varepsilon>0\),
\begin{equation}\label{eq:asymp-support-budget}
    \sum_{j\geq3}\varepsilon_j<\frac{\varepsilon}{2},
    \quad
    \frac1{\varepsilon_j}\leq C_\varepsilon 2^j,
    \quad
    \frac1{\varepsilon_j}
    \leq C_\varepsilon(1+|n|)^2,
    \quad n\in B_j.
\end{equation}

Let \(j\geq3\), and suppose that \(F\in L^2(\R)\) has bounded support and satisfies
\[
    \widehat F(\lambda_n)=0,
    \quad n\in B_3\cup\cdots\cup B_{j-1}.
\]
We construct a correction \(g_j\) which adds zeros on \(B_j\), preserves all the preceding zeros, and does not affect the blocks \(B_{j+1},\ldots,B_{4j}\). If \(B_j=\varnothing\), we simply take \(g_j=0\).

For \(n\in B_j\), we have
\[
    \frac14\leq2^{j-1}|\delta_n|<\frac12.
\]
Since \(\lambda_n=n+\delta_n\), it follows that
\begin{equation}\label{eq:asymp-dyadic-multiplier}
\begin{aligned}
    c
    &\leq
    \left|e^{2\pi i2^{j-1}\lambda_n}-1\right|
    \leq C,
    && n\in B_j,\\
    \left|e^{2\pi i2^{j-1}\lambda_n}-1\right|
    &\leq C2^j|\delta_n|,
    && n\in\Z.
\end{aligned}
\end{equation}
The first estimate allows us to divide by the multiplier on \(B_j\).
The second will control the effect of the correction on later blocks.

Choose an integer \(M_j\) sufficiently large that both intervals
\[
    [M_j,M_j+1]
    \quad\text{and}\quad
    [M_j-2^{j-1},M_j-2^{j-1}+1]
\]
are disjoint from \(\supp (F)\). Since \(j\ge3\), these intervals
are also disjoint from each other.
Applying Lemma~\ref{lem:asymp-thin-interpolation}, followed by a
translation to \([M_j,M_j+1]\), we obtain a function \(h\), supported
on a set of measure at most \(\varepsilon_j/2\), such that
\[
    \widehat h(\lambda_n)=
    \begin{cases}
        -\widehat F(\lambda_n)/
        \bigl(e^{2\pi i2^{j-1}\lambda_n}-1\bigr),
        &n\in B_j,\\[1mm]
        0,
        &n\in(B_3\cup\cdots\cup B_{4j})\setminus B_j,
    \end{cases}
\]
and
\[
    \norm{h}_2^2
    \leq\frac{C}{\varepsilon_j}
    \sum_{n\in B_j}|\widehat F(\lambda_n)|^2.
\]
Here the prescribed values are modified by the appropriate unimodular
factors before translating, so that the displayed identities hold for
the translated function.

Now consider
\[
    g_j(t)=h(t+2^{j-1})-h(t),
\]
which satisfies
\[
    \widehat g_j(x)
    =\bigl(e^{2\pi i2^{j-1}x}-1\bigr)\widehat h(x).
\]
Consequently, \(\widehat{F+g_j}\) vanishes on
\(B_3\cup\cdots\cup B_j\), while \(\widehat g_j\) vanishes on
\(B_{j+1}\cup\cdots\cup B_{4j}\). Since the two translates of \(h\)
have disjoint supports,
\begin{equation}\label{eq:asymp-one-block}
    |\supp (g_j)|\leq\varepsilon_j,
    \quad
    \norm{g_j}_2^2
    \leq\frac{C}{\varepsilon_j}
    \sum_{n\in B_j}|\widehat F(\lambda_n)|^2.
\end{equation}

It remains to estimate the effect of \(g_j\) on the more distant
blocks. If \(k>4j\) and \(n\in B_k\), then
\eqref{eq:asymp-support-budget} and
\eqref{eq:asymp-dyadic-multiplier} give
\[
    \frac{
        |e^{2\pi i2^{j-1}\lambda_n}-1|^2
    }{\varepsilon_k}
    \leq C_\varepsilon4^{-j}.
\]
Using the Bessel estimate \eqref{eq:asymp-bessel}, after translating
\(h\) back to an interval of length one, we obtain
\begin{equation}\label{eq:asymp-far-blocks}
\begin{aligned}
    \sum_{k>4j}\frac1{\varepsilon_k}
    \sum_{n\in B_k}|\widehat g_j(\lambda_n)|^2
    &=
    \sum_{k>4j}\frac1{\varepsilon_k}
    \sum_{n\in B_k}
    |e^{2\pi i2^{j-1}\lambda_n}-1|^2
    |\widehat h(\lambda_n)|^2\\
    &\leq
    C_\varepsilon4^{-j}
    \sum_{k>4j}\sum_{n\in B_k}
    |\widehat h(\lambda_n)|^2\\
    &\leq
    \frac{C_\varepsilon4^{-j}}{\varepsilon_j}
    \sum_{n\in B_j}|\widehat F(\lambda_n)|^2.
\end{aligned}
\end{equation}
This ensures that the corrections are square summable.
\medskip

\noindent\textbf{Step 2: Corrections on every block.}
Choose \(J\geq3\) so large that
\[
    C_\varepsilon^2\sum_{j\geq J}4^{-j}\leq\frac14,
\]
where \(C_\varepsilon\) is large enough for the estimates above. Take
a nonzero function \(\psi\in C_c^\infty(\R)\) with
\[
    |\supp(\psi)|<\frac{\varepsilon}{2},
\]
and define
\[
    F_{J-1}
    =
    \prod_{n\in B_3\cup\cdots\cup B_{J-1}}
    \left(\frac1{2\pi i}\frac{d}{dt}-\lambda_n\right)\psi.
\]
Then \(F_{J-1}\neq0\), its support is contained in \(\supp(\psi)\), and
its Fourier transform vanishes on
\(B_3\cup\cdots\cup B_{J-1}\). Because \(\psi\) is smooth, these
finitely many zeros can be introduced by differentiation without
enlarging its support. The later corrections are only \(L^2\), which
is why we use the interpolation construction from Step~1 for the
remaining blocks.

Since \(\widehat F_{J-1}\) is rapidly decreasing, the last estimate in
\eqref{eq:asymp-support-budget} gives
\[
    \sum_{j\geq J}\frac1{\varepsilon_j}
    \sum_{n\in B_j}
    |\widehat F_{J-1}(\lambda_n)|^2<\infty.
\]
Multiplying \(F_{J-1}\) by a nonzero constant, we may assume that
\begin{equation}\label{eq:asymp-initial-energy}
    \sum_{j\geq J}\frac1{\varepsilon_j}
    \sum_{n\in B_j}
    |\widehat F_{J-1}(\lambda_n)|^2\leq1.
\end{equation}

We now proceed through the blocks \(B_J,B_{J+1},\ldots\). At stage
\(j\), apply Step~1 to \(F_{j-1}\), place the correction \(g_j\) on a
part of the real line disjoint from all the preceding supports, and set
\[
    F_j=F_{j-1}+g_j.
\]
Thus each step adds one new block of zeros, preserves all preceding
zeros, and increases the measure of the support by at most
\(\varepsilon_j\).

We next verify that the corrections are square summable. For \(L\geq J\),
set
\[
    A_L=
    \left(
        \sum_{j=J}^{L}\frac1{\varepsilon_j}
        \sum_{n\in B_j}
        |\widehat F_{j-1}(\lambda_n)|^2
    \right)^{1/2}.
\]
For \(t<j\), the correction \(g_t\) vanishes on \(B_j\) whenever
\(j\leq4t\). Thus only corrections with \(j>4t\) contribute to the
values on \(B_j\). Minkowski's inequality,
\eqref{eq:asymp-initial-energy}, and
\eqref{eq:asymp-far-blocks} therefore give
\[
    A_L
    \leq
    1+
    \sum_{t=J}^{L-1}
    \frac{C_\varepsilon2^{-t}}{\sqrt{\varepsilon_t}}
    \left(
        \sum_{n\in B_t}
        |\widehat F_{t-1}(\lambda_n)|^2
    \right)^{1/2}.
\]
By Cauchy-Schwarz and the choice of \(J\),
\[
    A_L
    \leq
    1+
    C_\varepsilon
    \left(\sum_{t\geq J}4^{-t}\right)^{1/2}A_L
    \leq1+\frac12A_L.
\]
Hence \(A_L\leq2\), uniformly in \(L\), and consequently
\[
    \sum_{j\geq J}\frac1{\varepsilon_j}
    \sum_{n\in B_j}
    |\widehat F_{j-1}(\lambda_n)|^2
    \leq4.
\]
It now follows from \eqref{eq:asymp-one-block} that
\[
    \sum_{j\geq J}\norm{g_j}_2^2<\infty.
\]

\medskip
\noindent\textbf{Step 3: Passing to the limit.}
The supports of the corrections were chosen pairwise disjoint and
disjoint from \(\supp (F_{J-1})\). Therefore, we have that
\[
    F=F_{J-1}+\sum_{j\geq J}g_j
\]
converges in \(L^2(\R)\). Since \(F\) agrees with
the nonzero function \(F_{J-1}\) on \(\supp (F_{J-1})\), it follows that $F \neq 0$.

In addition, we observe that
\[
    \sum_{j\geq J}\norm{g_j}_1
    \leq
    \left(\sum_{j\geq J}|\supp (g_j)|\right)^{1/2}
    \left(\sum_{j\geq J}\norm{g_j}_2^2\right)^{1/2}
    <\infty.
\]
It follows that \(\widehat F_j\) converges uniformly to \(\widehat F\).
At stage \(j\), the Fourier transform of \(F_j\) vanishes on
\(B_3\cup\cdots\cup B_j\), and every subsequent correction preserves
these zeros. Since the blocks \(B_j\), where \(j\geq3\), cover \(\Z\), we
conclude that for all $n \in \Z$ it holds that $\widehat F(\lambda_n)=0$.

Finally,
\[
    |\supp (F)|
    \leq
    |\supp (F_{J-1})|+\sum_{j\geq J}\varepsilon_j
    <\varepsilon.
\]
Taking \(S=\supp (F)\), the nonzero function \(F\in L^2(S)\) is
orthogonal to \(E(\Lambda)\), and the theorem follows.
\end{proof}

\end{document}
